\documentclass[a4paper,12pt]{article}

% ------------------------------------------------
% Кодировка, язык, математика
% ------------------------------------------------
\usepackage[T2A]{fontenc}
\usepackage[utf8]{inputenc}
\usepackage[russian]{babel}
\usepackage{amsmath,amssymb,amsthm,mathtools}
\usepackage{icomma}

% ------------------------------------------------
% Шрифт и типографика
% ------------------------------------------------
\usepackage{helvet}
\renewcommand{\familydefault}{\sfdefault}
\usepackage{microtype}
\usepackage{setspace}
\onehalfspacing
\usepackage{parskip}

% ------------------------------------------------
% Страница
% ------------------------------------------------
\usepackage[
  a4paper,
  left=20mm,
  right=20mm,
  top=20mm,
  bottom=22mm
]{geometry}

% ------------------------------------------------
% Цвета и графика
% ------------------------------------------------
\usepackage{xcolor}
\definecolor{accent}{HTML}{008080}
\definecolor{darkaccent}{HTML}{004D4D}
\definecolor{textgray}{HTML}{333333}
\definecolor{softgray}{HTML}{777777}
\definecolor{lightaccent}{HTML}{E7F2F2}

\usepackage{tikz}
\usetikzlibrary{calc}

% Подключены согласно техническому стандарту.
% В данном пособии специальные геометрические чертежи
% и графики функций не требуются.
\usepackage{tkz-euclide}
\usepackage{pgfplots}
\pgfplotsset{compat=1.18}

% ------------------------------------------------
% Таблицы и списки
% ------------------------------------------------
\usepackage{tabularx,booktabs,longtable}
\usepackage{enumitem}

\setlist[itemize]{
  leftmargin=6mm,
  itemsep=3pt,
  topsep=4pt
}

\setlist[enumerate]{
  leftmargin=8mm,
  itemsep=7pt,
  topsep=5pt
}

% ------------------------------------------------
% Заголовки и колонтитулы
% ------------------------------------------------
\usepackage{fancyhdr}
\usepackage{hyperref}

\hypersetup{
  colorlinks=true,
  linkcolor=darkaccent,
  urlcolor=darkaccent
}

\pagestyle{fancy}
\fancyhf{}
\fancyhead[L]{\small\textcolor{softgray}{Ксения Васильченко}}
\fancyhead[R]{\small\textcolor{softgray}{Математика, 6 класс}}
\fancyfoot[C]{\small\textcolor{softgray}{\thepage}}
\renewcommand{\headrulewidth}{0.3pt}
\renewcommand{\headrule}{%
  \hbox to\headwidth{%
    \color{accent}\leaders\hrule height \headrulewidth\hfill}}

% ------------------------------------------------
% Собственные команды
% ------------------------------------------------
\newcommand{\topic}[1]{%
  \vspace{0.7em}
  {\Large\bfseries\color{darkaccent}#1}\par
  \vspace{0.25em}
  {\color{accent}\rule{\linewidth}{0.5pt}}
  \vspace{0.55em}
}

\newcommand{\subtopic}[1]{%
  \vspace{0.4em}
  {\large\bfseries\color{darkaccent}#1}\par
  \vspace{0.15em}
}

\newcommand{\important}[1]{%
  \par
  \vspace{0.25em}
  \noindent
  \textcolor{darkaccent}{\textbf{Важно.}} #1
  \par
}

\newcommand{\exampletitle}{%
  \par
  \vspace{0.3em}
  \noindent
  \textcolor{accent}{\textbf{Пример.}}\enspace
}

\newcommand{\mistaketitle}{%
  \par
  \vspace{0.3em}
  \noindent
  \textcolor{darkaccent}{\textbf{Типичная ошибка.}}\enspace
}

\newcommand{\checkitem}{\(\square\)}

\begin{document}

% ================================================================
% ТИТУЛЬНЫЙ ЛИСТ
% ================================================================

\begin{titlepage}
\thispagestyle{empty}

\begin{center}

\vspace*{22mm}

{\Large\textcolor{accent}{МАТЕМАТИКА · 6 КЛАСС}}

\vspace{15mm}

{\fontsize{31}{37}\selectfont
\bfseries\textcolor{darkaccent}{Чек-лист}}

\vspace{7mm}

{\LARGE
Дроби, среднее арифметическое,\\[2mm]
простейшие уравнения и средняя скорость}

\vspace{15mm}

{\large
Повторение базовых вычислительных навыков\\
и подготовка к следующим темам курса
}

\vfill

{\large
\textbf{Лёвин Артём Александрович}\\[2mm]
эксклюзивно для Ксении Васильченко
}

\vspace{9mm}

{\normalsize \today}

\vspace{25mm}

\end{center}

% Кривая Безье в нижней части титульного листа
\begin{tikzpicture}[remember picture,overlay]
\draw[
  accent,
  line width=1.1pt
]
([xshift=18mm,yshift=17mm]current page.south west)
.. controls
([xshift=55mm,yshift=29mm]current page.south west)
and
([xshift=-65mm,yshift=7mm]current page.south east)
..
([xshift=-18mm,yshift=19mm]current page.south east);
\end{tikzpicture}

\end{titlepage}

% ================================================================
% ОСНОВНОЙ ЧЕК-ЛИСТ
% ================================================================

\section*{\textcolor{darkaccent}{Главное после занятия}}

После повторения Вы должны уметь:

\begin{itemize}
  \item[\checkitem] выполнять действия с десятичными дробями;
  \item[\checkitem] читать разряды десятичной дроби;
  \item[\checkitem] складывать и вычитать обыкновенные дроби;
  \item[\checkitem] сокращать дроби;
  \item[\checkitem] переводить неправильную дробь в смешанное число;
  \item[\checkitem] умножать и делить обыкновенные дроби;
  \item[\checkitem] находить среднее арифметическое;
  \item[\checkitem] находить неизвестное число по известному среднему арифметическому;
  \item[\checkitem] переводить словесные отношения между величинами на математический язык;
  \item[\checkitem] решать простейшие уравнения;
  \item[\checkitem] находить среднюю скорость по всему пути;
  \item[\checkitem] переводить минуты в часы и метры в километры.
\end{itemize}

% ----------------------------------------------------------------
\topic{1. Десятичные дроби}

\subtopic{1.1. Сложение и вычитание}

При сложении и вычитании десятичных дробей числа записываем
\textbf{запятая под запятой}.

При необходимости справа дописываем нули:

\[
1{,}2=1{,}20.
\]

\exampletitle
\[
1{,}57+1{,}20=2{,}77.
\]

\[
1{,}57-1{,}20=0{,}37.
\]

\important{
Количество цифр после запятой можно временно сделать одинаковым,
дописав справа нули. Значение числа при этом не изменяется.
}

\subtopic{1.2. Названия разрядов}

\begin{center}
\renewcommand{\arraystretch}{1.3}
\begin{tabularx}{0.84\linewidth}{@{}cX@{}}
\toprule
\textbf{Запись} & \textbf{Дробная часть}\\
\midrule
$1{,}2$ & две десятых\\
$1{,}23$ & двадцать три сотых\\
$1{,}234$ & двести тридцать четыре тысячных\\
$1{,}2345$ & две тысячи триста сорок пять десятитысячных\\
\bottomrule
\end{tabularx}
\end{center}

Число цифр после запятой показывает разряд последней цифры.

\subtopic{1.3. Умножение десятичных дробей}

Алгоритм:

\begin{enumerate}
  \item временно не обращаем внимания на запятые;
  \item выполняем обычное умножение;
  \item считаем общее количество цифр после запятых в множителях;
  \item отделяем столько же цифр справа в ответе.
\end{enumerate}

\exampletitle
\[
1{,}2\cdot1{,}57.
\]

В исходных множителях всего
\[
1+2=3
\]
цифры после запятых, поэтому в результате также должно быть
три цифры после запятой:

\[
1{,}2\cdot1{,}57=1{,}884.
\]

\subtopic{1.4. Деление десятичного числа на натуральное}

\exampletitle
\[
1{,}2:4=0{,}3.
\]

Проверка:

\[
0{,}3\cdot4=1{,}2.
\]

% ----------------------------------------------------------------
\topic{2. Обыкновенные дроби}

\subtopic{2.1. Сложение и вычитание}

Для сложения и вычитания дробей знаменатели должны быть одинаковыми.

Например,

\[
\frac12+\frac13
=
\frac{3}{6}+\frac{2}{6}
=
\frac56.
\]

И аналогично:

\[
\frac12-\frac13
=
\frac36-\frac26
=
\frac16.
\]

\exampletitle

\[
\frac14-\frac15
=
\frac5{20}-\frac4{20}
=
\frac1{20}.
\]

\important{
При сложении и вычитании дробей числители нельзя складывать или
вычитать до приведения дробей к общему знаменателю.
}

\subtopic{2.2. Целое число как дробь}

Любое целое число можно записать со знаменателем~$1$:

\[
2=\frac21,\qquad
5=\frac51.
\]

Поэтому

\[
2-\frac13
=
\frac21-\frac13
=
\frac63-\frac13
=
\frac53.
\]

% ----------------------------------------------------------------
\topic{3. Неправильная дробь и смешанное число}

Если числитель больше знаменателя или равен ему, дробь называется
\textbf{неправильной}.

Например,

\[
\frac53,\qquad \frac{10}{4}.
\]

Чтобы представить неправильную дробь в виде смешанного числа,
делим числитель на знаменатель с остатком.

\exampletitle

\[
5:3=1\text{ (ост. }2\text{)}.
\]

Следовательно,

\[
\frac53=1\frac23.
\]

Здесь:

\begin{itemize}
  \item $1$ --- целая часть;
  \item $2$ --- остаток, который становится новым числителем;
  \item $3$ --- прежний знаменатель.
\end{itemize}

Ещё один пример:

\[
\frac{10}{4}=2\frac24=2\frac12.
\]

\important{
После выделения целой части дробную часть желательно сократить,
если это возможно.
}

% ----------------------------------------------------------------
\topic{4. Сокращение дробей}

Чтобы сократить дробь, числитель и знаменатель делим на один и тот же
общий множитель.

\exampletitle

\[
\frac{8}{6}
=
\frac{2\cdot4}{2\cdot3}
=
\frac43.
\]

Ещё пример:

\[
\frac{12}{18}
=
\frac{6\cdot2}{6\cdot3}
=
\frac23.
\]

Можно сокращать постепенно:

\[
\frac{12}{18}
=
\frac69
=
\frac23.
\]

\mistaketitle
Сократить можно только \textbf{множители}.
В сумме или разности отдельные числа ``зачёркивать'' нельзя.

% ----------------------------------------------------------------
\topic{5. Умножение и деление обыкновенных дробей}

\subtopic{5.1. Умножение}

\[
\frac ab\cdot\frac cd
=
\frac{ac}{bd}.
\]

\exampletitle

\[
\frac12\cdot\frac35
=
\frac{3}{10}.
\]

Перед умножением полезно проверить, можно ли выполнить сокращение.

\subtopic{5.2. Деление}

При делении на дробь заменяем деление умножением на
\textbf{обратную дробь}:

\[
\frac ab:\frac cd
=
\frac ab\cdot\frac dc,
\qquad c\ne0,\ d\ne0.
\]

\exampletitle

\[
\frac12:\frac35
=
\frac12\cdot\frac53
=
\frac56.
\]

\important{
Переворачивается только \textbf{вторая} дробь.
}

% ----------------------------------------------------------------
\topic{6. Среднее арифметическое}

Среднее арифметическое нескольких чисел равно их сумме,
делённой на количество чисел:

\[
\boxed{
\text{Среднее арифметическое}
=
\frac{\text{сумма всех чисел}}
     {\text{количество чисел}}
}
\]

\exampletitle

Для чисел
\[
1,\ 2,\ 3,\ 4,\ 5
\]
получаем:

\[
\frac{1+2+3+4+5}{5}
=
\frac{15}{5}
=
3.
\]

\subtopic{Полезный вычислительный приём}

При сложении нескольких чисел ищите удобные пары.

Например,

\[
4{,}9+5{,}1+5+4{,}8+5{,}2.
\]

Удобно сгруппировать:

\[
(4{,}9+5{,}1)+(4{,}8+5{,}2)+5
=
10+10+5
=
25.
\]

Следовательно,

\[
\frac{25}{5}=5.
\]

Такой способ уменьшает количество вычислений и снижает вероятность ошибки.

% ----------------------------------------------------------------
\topic{7. Как найти неизвестное число по среднему арифметическому}

Пусть одно число равно $6{,}4$, второе равно $x$, а их среднее
арифметическое равно $3{,}25$.

Составляем уравнение:

\[
\frac{6{,}4+x}{2}=3{,}25.
\]

Умножаем обе части на $2$:

\[
6{,}4+x=6{,}5.
\]

Следовательно,

\[
x=6{,}5-6{,}4=0{,}1.
\]

Ответ:

\[
\boxed{x=0{,}1}.
\]

\important{
Сначала переводите условие задачи на математический язык.
Неизвестную величину удобно обозначать буквой $x$.
}

\subtopic{Ещё один пример}

Одно число равно $5$, другое неизвестно, а среднее арифметическое
равно $3{,}5$:

\[
\frac{5+x}{2}=3{,}5.
\]

Тогда

\[
5+x=7,
\]

\[
x=2.
\]

% ----------------------------------------------------------------
\topic{8. Базовые связи между компонентами действий}

Эти правила особенно важны при решении уравнений.

\subtopic{8.1. Сложение}

Если

\[
a+b=c,
\]

то

\[
a=c-b,
\qquad
b=c-a.
\]

\subtopic{8.2. Вычитание}

Если

\[
a-b=c,
\]

то

\[
a=b+c,
\qquad
b=a-c.
\]

\subtopic{8.3. Умножение}

Если

\[
ab=c
\]

и $a\ne0$, то

\[
b=\frac ca.
\]

\subtopic{8.4. Деление}

Если

\[
\frac ab=c,
\qquad b\ne0,
\]

то

\[
a=bc.
\]

Если дополнительно $c\ne0$, то

\[
b=\frac ac.
\]

\important{
В дроби знаменатель никогда не может быть равен нулю.
}

% ----------------------------------------------------------------
\topic{9. Простейшие уравнения}

Рассмотрим несколько базовых моделей.

\[
x-10=50
\]

\[
x=50+10=60.
\]

\vspace{0.4em}

\[
\frac{x}{5}=10
\]

\[
x=5\cdot10=50.
\]

\vspace{0.4em}

\[
3x=9
\]

\[
x=9:3=3.
\]

\vspace{0.4em}

Более составной пример:

\[
\frac{x-10}{3}=50.
\]

Сначала находим весь числитель:

\[
x-10=150.
\]

Затем:

\[
x=160.
\]

\important{
Составное выражение удобно временно воспринимать как один объект.
Например, в дроби
\[
\frac{x-10}{3}=50
\]
весь числитель $x-10$ играет роль делимого.
}

% ----------------------------------------------------------------
\topic{10. Перевод текста задачи на математический язык}

Пусть величины обозначены буквами $k$ и $t$.

\begin{center}
\renewcommand{\arraystretch}{1.4}
\begin{tabularx}{0.9\linewidth}{@{}Xc@{}}
\toprule
\textbf{Фраза} & \textbf{Математическая запись}\\
\midrule
$k$ больше $t$ на $50$ & $k=t+50$\\
$k$ меньше $t$ на $50$ & $k=t-50$\\
$k$ больше $t$ в $50$ раз & $k=50t$\\
$k$ меньше $t$ в $50$ раз & $k=\dfrac{t}{50}$\\
\bottomrule
\end{tabularx}
\end{center}

\important{
Фразы \textbf{«на столько-то»} и \textbf{«во столько-то раз»}
описывают разные действия:
\[
\text{«на»} \longrightarrow +\ \text{или}\ -,
\]
\[
\text{«в раз»} \longrightarrow \cdot\ \text{или}\ :.
\]
}

% ----------------------------------------------------------------
\topic{11. Два неизвестных числа и связь между ними}

Задача:

\begin{quote}
Среднее арифметическое двух чисел равно $146$.
Одно число больше другого на $22$.
Найдите эти числа.
\end{quote}

Пусть меньшее число равно $y$.

Тогда большее:

\[
y+22.
\]

По определению среднего арифметического:

\[
\frac{y+(y+22)}{2}=146.
\]

Упрощаем:

\[
\frac{2y+22}{2}=146.
\]

Умножаем на $2$:

\[
2y+22=292.
\]

\[
2y=270.
\]

\[
y=135.
\]

Тогда второе число:

\[
135+22=157.
\]

Проверка:

\[
\frac{135+157}{2}
=
\frac{292}{2}
=
146.
\]

Ответ:

\[
\boxed{135\text{ и }157}.
\]

\important{
После решения текстовой задачи полезно сделать проверку по исходному условию.
}

% ----------------------------------------------------------------
\topic{12. Средняя скорость}

Средняя скорость за весь путь находится по формуле

\[
\boxed{
v_{\text{ср}}
=
\frac{S_{\text{общ}}}{t_{\text{общ}}}
}
\]

где

\[
S_{\text{общ}}
=
S_1+S_2+\ldots,
\]

\[
t_{\text{общ}}
=
t_1+t_2+\ldots.
\]

То есть

\[
\boxed{
v_{\text{ср}}
=
\frac{S_1+S_2+\cdots+S_n}
     {t_1+t_2+\cdots+t_n}
}
\]

\important{
Среднюю скорость обычно нельзя находить как обычное среднее
от скоростей отдельных участков.
Нужно разделить \textbf{весь путь} на \textbf{всё время}.
}

\subtopic{Пример с переводом минут в часы}

Велосипедист проехал:

\[
1{,}2\text{ км за }6\text{ мин},
\]

\[
5{,}3\text{ км за }12\text{ мин},
\]

\[
2{,}3\text{ км за }15\text{ мин}.
\]

Общий путь:

\[
S_{\text{общ}}
=
1{,}2+5{,}3+2{,}3
=
8{,}8\text{ км}.
\]

Общее время:

\[
t_{\text{общ}}
=
6+12+15
=
33\text{ мин}.
\]

Для ответа в километрах в час переводим минуты в часы:

\[
33\text{ мин}
=
\frac{33}{60}\text{ ч}
=
\frac{11}{20}\text{ ч}.
\]

Следовательно,

\[
v_{\text{ср}}
=
8{,}8:\frac{11}{20}.
\]

Заменяем деление умножением на обратную дробь:

\[
v_{\text{ср}}
=
8{,}8\cdot\frac{20}{11}.
\]

Так как

\[
8{,}8=\frac{44}{5},
\]

получаем:

\[
v_{\text{ср}}
=
\frac{44}{5}\cdot\frac{20}{11}
=
16.
\]

Ответ:

\[
\boxed{16\text{ км/ч}}.
\]

% ----------------------------------------------------------------
\topic{13. Перевод единиц}

\subtopic{Минуты в часы}

В одном часе $60$ минут, поэтому

\[
\boxed{
t_{\text{ч}}=\frac{t_{\text{мин}}}{60}
}
\]

Например,

\[
30\text{ мин}
=
\frac{30}{60}\text{ ч}
=
\frac12\text{ ч}.
\]

\[
45\text{ мин}
=
\frac{45}{60}\text{ ч}
=
\frac34\text{ ч}.
\]

\subtopic{Метры в километры}

В одном километре $1000$ метров:

\[
\boxed{
S_{\text{км}}=\frac{S_{\text{м}}}{1000}
}
\]

Например,

\[
1500\text{ м}
=
\frac{1500}{1000}\text{ км}
=
1{,}5\text{ км}.
\]

% ----------------------------------------------------------------
\topic{14. Что особенно важно запомнить}

\begin{enumerate}
  \item В десятичных дробях при сложении и вычитании:
  \[
  \text{запятая под запятой}.
  \]

  \item Среднее арифметическое:
  \[
  \frac{\text{сумма чисел}}{\text{количество чисел}}.
  \]

  \item При делении на дробь вторую дробь переворачиваем:
  \[
  \frac ab:\frac cd
  =
  \frac ab\cdot\frac dc.
  \]

  \item Если величина больше другой \textbf{на} $a$:
  \[
  x=y+a.
  \]

  \item Если величина больше другой \textbf{в} $a$ раз:
  \[
  x=ay.
  \]

  \item Средняя скорость:
  \[
  v_{\text{ср}}
  =
  \frac{\text{весь путь}}{\text{всё время}}.
  \]

  \item Минуты переводим в часы делением на $60$.

  \item Метры переводим в километры делением на $1000$.

  \item Все промежуточные вычисления желательно записывать:
  это помогает заметить арифметическую ошибку.
\end{enumerate}

% ================================================================
% ТРЕНИРОВОЧНЫЙ БЛОК
% ================================================================

\clearpage

\section*{\textcolor{darkaccent}{Тренировочный блок}}

{\color{accent}\rule{\linewidth}{0.6pt}}

\vspace{0.5em}

Выполните задания последовательно.
Решения записывайте подробно, включая промежуточные вычисления.

\begin{enumerate}[label=\textbf{\arabic*.}]

  \item
  Выполните действия:
  \[
  3{,}47+2{,}8,
  \qquad
  5{,}2-1{,}76,
  \]
  \[
  \frac23-\frac14,
  \qquad
  \frac{18}{24}.
  \]
  Последнюю дробь сократите.

  \item
  Найдите среднее арифметическое чисел
  \[
  6{,}8;\quad
  7{,}2;\quad
  5{,}9;\quad
  6{,}1;\quad
  7.
  \]
  Постарайтесь сначала подобрать удобные пары для сложения.

  \item
  Одно из двух чисел равно $7{,}4$, а их среднее арифметическое
  равно $5{,}2$.
  Найдите второе число.

  \item
  Среднее арифметическое двух чисел равно $84$.
  Одно число больше другого на $18$.
  Найдите оба числа.

  \item
  Велосипедист проехал первый участок длиной $2{,}4$ км за
  $8$ минут, второй участок длиной $3{,}6$ км за $12$ минут,
  а третий участок длиной $2$ км за $10$ минут.

  Найдите среднюю скорость велосипедиста на всём пути.
  Ответ выразите в километрах в час.

\end{enumerate}

\vfill

\begin{center}
\textcolor{softgray}{
Перед проверкой убедитесь, что в задачах на движение
все единицы измерения согласованы.
}
\end{center}

% ================================================================
% ОТВЕТЫ
% ================================================================

\clearpage

\section*{\textcolor{darkaccent}{Ответы к тренировочному блоку}}

{\color{accent}\rule{\linewidth}{0.6pt}}

\vspace{1em}

\begin{table}[ht]
\centering
\caption{Ответы}
\renewcommand{\arraystretch}{1.55}
\begin{tabularx}{\linewidth}{@{}cX@{}}
\toprule
\textbf{№} & \textbf{Ответ}\\
\midrule

1 &
\[
3{,}47+2{,}8=6{,}27;
\qquad
5{,}2-1{,}76=3{,}44;
\]
\[
\frac23-\frac14=\frac5{12};
\qquad
\frac{18}{24}=\frac34.
\]
\\

2 &
\[
6{,}6.
\]
\\

3 &
\[
3.
\]
\\

4 &
\[
75\text{ и }93.
\]
\\

5 &
Общий путь:
\[
2{,}4+3{,}6+2=8\text{ км}.
\]

Общее время:
\[
8+12+10=30\text{ мин}
=
\frac12\text{ ч}.
\]

Средняя скорость:
\[
v_{\text{ср}}
=
8:\frac12
=
16\text{ км/ч}.
\]

\[
\boxed{16\text{ км/ч}}
\]
\\

\bottomrule
\end{tabularx}
\end{table}

\vfill


\end{document}